% ============================================================
%  Bab 2 – Dasar Teori
% ============================================================

\section{Dasar Teori}

\subsection{Supervised Learning}
% TODO: Definisi, paradigma, hubungan dengan klasifikasi.

\subsection{Klasifikasi}
% TODO: Pengertian umum klasifikasi biner vs.\ multiclass.

\subsection{K-Nearest Neighbors}
% TODO: Deskripsi algoritma, perhitungan jarak, pemilihan k.

\subsection{Decision Tree Learning}

\subsubsection{Iterative Dichotomiser 3 (ID3)}
% TODO: Kriteria information gain, entropy.

\subsubsection{C4.5}
% TODO: Gain ratio, penanganan nilai kontinu dan hilang.

\subsubsection{Classification and Regression Tree (CART)}
% TODO: Gini impurity, pruning.

\subsection{Logistic Regression}

\subsubsection{Regularisasi L1}
\subsubsection{Regularisasi L2}
\subsubsection{Elastic Net}

\subsection{Support Vector Machine}

\subsubsection{Kernel Linear}
\subsubsection{Kernel Polinomial}
\subsubsection{Kernel Gaussian (RBF)}

\subsection{Bayesian Network (Naive Bayes)}
% TODO: Teorema Bayes, asumsi conditional independence, Laplace smoothing.

\subsection{Validasi dan Metrik Evaluasi}
% TODO: Train-test split, k-fold CV, accuracy, precision, recall, F1, ROC-AUC.

